% 分野別類題: 複素指数関数の積分・sinc関数 解答例
% コンパイル: TEXINPUTS="../../../00_common:$TEXINPUTS" latexmk -xelatex answers.tex
\documentclass[a4paper,11pt]{article}
\usepackage{preamble}
\usepackage{shoakubox}

\begin{document}

\kamokumei{類題演習 解答例 --- Complex Exponential Integrals and the Sinc Function}

\daimon{【解法】複素指数関数の積分と sinc 関数の導出}

\noindent
オイラーの公式
\[
e^{i\theta} = \cos\theta + i\sin\theta
\]
を用いると、指数関数の積分公式 $\displaystyle\int e^{\lambda x}\,dx = \frac{1}{\lambda}e^{\lambda x} + C$（$\lambda$ は複素定数、$\lambda \neq 0$）は $\lambda$ が純虚数や複素数であっても形式的にそのまま成り立つ。実際、$\lambda = a + ib$ とすると
\[
\frac{d}{dx}e^{\lambda x} = \frac{d}{dx}\left(e^{ax}\cos bx + i\,e^{ax}\sin bx\right)
= \lambda\,e^{\lambda x}
\]
が実部・虚部それぞれの微分公式から確かめられるため、複素指数関数についても通常の指数関数と同じ積分公式・微分公式が使える。

\medskip
\noindent
\textbf{(1) 有限区間の複素指数積分と sinc 関数。}
$k \neq 0$ のとき $\lambda = -ik$ として
\[
\int_{-a}^{a} e^{-ikx}\,dx = \left[\frac{e^{-ikx}}{-ik}\right]_{-a}^{a}
= \frac{e^{-ika} - e^{ika}}{-ik}
= \frac{-2i\sin(ka)}{-ik}
= \frac{2\sin(ka)}{k}
\]
ここでオイラーの公式より $e^{-ika} - e^{ika} = -2i\sin(ka)$ を用いた。これを
\[
\frac{2\sin(ka)}{k} = 2a\cdot\frac{\sin(ka)}{ka} = 2a\,\mathrm{sinc}(ka)
\]
と変形すれば、有限区間 $[-a,a]$ 上の複素指数関数の積分が $\mathrm{sinc}$ 関数で表されることがわかる。これは矩形パルスのフーリエ変換（問3）の核心部分である。

\medskip
\noindent
\textbf{(2) 実指数関数と三角関数の積の積分。}
$\displaystyle\int e^{ax}\cos(bx)\,dx$ を求めるには、$e^{ax}\cos(bx) = \mathrm{Re}\left(e^{(a+ib)x}\right)$ に注目し、複素指数関数
\[
\int e^{(a+ib)x}\,dx = \frac{1}{a+ib}e^{(a+ib)x} + C
\]
を計算してから実部を取ればよい。$\dfrac{1}{a+ib} = \dfrac{a-ib}{a^{2}+b^{2}}$ であることと $e^{(a+ib)x} = e^{ax}(\cos bx + i\sin bx)$ より、実部・虚部を比較すれば $\cos bx$, $\sin bx$ 双方の積分公式が同時に得られる。

\bigskip

\clearpage
\begin{mondaiboxS}{問1}
実数 $k \neq 0$ と正の定数 $a$ に対して、$\displaystyle\int_{-a}^{a} e^{-ikx}\,dx$ を計算し、その結果が $2a\,\mathrm{sinc}(ka)$ と表せることを示せ。
\end{mondaiboxS}
\kaitou
$\lambda = -ik \neq 0$ として原始関数を求めると
\[
\int_{-a}^{a} e^{-ikx}\,dx
= \left[\frac{e^{-ikx}}{-ik}\right]_{-a}^{a}
= \frac{e^{-ika} - e^{ika}}{-ik}
\]
オイラーの公式より $e^{-ika} = \cos(ka) - i\sin(ka)$, $e^{ika} = \cos(ka) + i\sin(ka)$ であるから
\[
e^{-ika} - e^{ika} = -2i\sin(ka)
\]
よって
\[
\int_{-a}^{a} e^{-ikx}\,dx = \frac{-2i\sin(ka)}{-ik} = \frac{2\sin(ka)}{k}
\]
これを $2a$ で括り直すと
\[
\frac{2\sin(ka)}{k} = 2a\cdot\frac{\sin(ka)}{ka}
\]
したがって
\[
\boxed{\; \int_{-a}^{a} e^{-ikx}\,dx = 2a\,\mathrm{sinc}(ka) \;}
\]
（検算: 被積分関数の実部 $\cos(kx)$ は偶関数なので $\displaystyle\int_{-a}^{a}\cos(kx)\,dx = 2\int_{0}^{a}\cos(kx)\,dx = \dfrac{2\sin(ka)}{k}$ となり実部が一致し、虚部 $-\sin(kx)$ は奇関数なので積分は $0$ となって結果が実数であることとも整合する。また $k \to 0$ の極限では $\mathrm{sinc}(0)=1$ より積分値は $2a$ に収束し、これは $\displaystyle\int_{-a}^{a}1\,dx = 2a$ と一致する。）

\clearpage
\begin{mondaiboxS}{問2}
実数 $a, b$（$a \neq 0$）に対して、$\displaystyle\int e^{ax}\cos(bx)\,dx$ を求めよ。
\end{mondaiboxS}
\kaitou
$e^{ax}\cos(bx) = \mathrm{Re}\left(e^{(a+ib)x}\right)$ であるから、まず複素指数関数の積分を計算する。$\lambda = a+ib \neq 0$ として
\[
\int e^{(a+ib)x}\,dx = \frac{1}{a+ib}e^{(a+ib)x} + C_{1}
\]
分母を実数化するため $\dfrac{1}{a+ib} = \dfrac{a-ib}{a^{2}+b^{2}}$ を用いると
\[
\frac{1}{a+ib}e^{(a+ib)x}
= \frac{a-ib}{a^{2}+b^{2}}\,e^{ax}\bigl(\cos bx + i\sin bx\bigr)
\]
右辺を展開して実部を取ると
\[
\mathrm{Re}\left[\frac{a-ib}{a^{2}+b^{2}}\,e^{ax}(\cos bx + i\sin bx)\right]
= \frac{e^{ax}}{a^{2}+b^{2}}\bigl(a\cos bx + b\sin bx\bigr)
\]
したがって
\[
\boxed{\; \int e^{ax}\cos(bx)\,dx = \frac{e^{ax}\bigl(a\cos bx + b\sin bx\bigr)}{a^{2}+b^{2}} + C \;}
\]
（検算: 右辺を $x$ で微分すると
\[
\frac{d}{dx}\left[\frac{e^{ax}(a\cos bx + b\sin bx)}{a^{2}+b^{2}}\right]
= \frac{a\,e^{ax}(a\cos bx + b\sin bx) + e^{ax}(-ab\sin bx + b^{2}\cos bx)}{a^{2}+b^{2}}
\]
分子を整理すると $e^{ax}\bigl[(a^{2}+b^{2})\cos bx + (ab - ab)\sin bx\bigr] = (a^{2}+b^{2})e^{ax}\cos bx$ となるので、全体は $e^{ax}\cos bx$ に一致し、確かに元の被積分関数に戻る。）

\clearpage
\begin{mondaiboxS}{問3}
高さ $1$、幅 $2a$ の矩形パルス関数 $f(x) = 1\;(|x|\le a)$, $f(x) = 0\;(|x|>a)$ のフーリエ変換 $F(k) = \displaystyle\int_{-\infty}^{\infty} f(x)\,e^{-ikx}\,dx$ を求め、$\mathrm{sinc}$ 関数で表せることを示せ。また $\displaystyle\lim_{k \to 0} F(k)$ を求めよ。
\end{mondaiboxS}
\kaitou
$f(x)$ は $|x|\le a$ でのみ $1$、それ以外で $0$ であるから
\[
F(k) = \int_{-\infty}^{\infty} f(x)\,e^{-ikx}\,dx = \int_{-a}^{a} e^{-ikx}\,dx
\]
これは問1で計算した積分そのものである。$k \neq 0$ のとき、問1の結果より
\[
F(k) = 2a\,\mathrm{sinc}(ka) = \frac{2\sin(ka)}{k}
\]
$k=0$ のときは
\[
F(0) = \int_{-a}^{a} 1\,dx = 2a
\]
であり、$\mathrm{sinc}(0)=1$ の定義とも整合するので、上式は $k=0$ を含めて全ての実数 $k$ で成り立つ。したがって
\[
\boxed{\; F(k) = 2a\,\mathrm{sinc}(ka) \;}
\]
また、$\mathrm{sinc}$ の極限 $\displaystyle\lim_{t\to 0}\frac{\sin t}{t} = 1$（またはロピタルの定理）を用いれば
\[
\lim_{k\to 0} F(k) = \lim_{k\to 0} 2a\,\mathrm{sinc}(ka) = 2a\cdot 1 = 2a
\]
すなわち
\[
\boxed{\; \lim_{k\to 0} F(k) = 2a \;}
\]
（検算: これは $F(0) = \int_{-\infty}^{\infty} f(x)\,dx = 2a$、すなわちパルスの「面積」と一致しており、フーリエ変換の定義から直接計算した値と極限操作による値が矛盾なく一致することが確認できる。）

\end{document}
